\documentclass[11pt]{article}

\usepackage[margin=1in]{geometry}
\usepackage{amsmath}
\usepackage{booktabs}
\usepackage{array}
\usepackage{listings}
\usepackage{xcolor}
\usepackage[hidelinks]{hyperref}

\lstdefinestyle{code}{
  basicstyle=\ttfamily\footnotesize,
  breaklines=true,
  columns=fullflexible,
  frame=single,
  numbers=left,
  numberstyle=\scriptsize\color{gray},
  showstringspaces=false,
  tabsize=4
}

\lstdefinestyle{output}{
  basicstyle=\ttfamily\scriptsize,
  breaklines=true,
  columns=fullflexible,
  frame=single,
  showstringspaces=false
}

\title{Computational Appendix\\
Less Left on the Shelf: Reducing Waste at Pasta \& Co.}
\author{Elisa Regueros \and Chandler Olson}
\date{November 2026}

\begin{document}
\maketitle

\section{Purpose}

This appendix records the data-processing rules, implementation, and exact
computer output supporting the preliminary results in the main report.  The
analysis uses an empirical 80\% service-level rule.  Sensitivity testing and
Monte Carlo simulation are planned for a later stage and are not included here.

\section{Data Construction and Checks}

The study period is September 13 through October 13, 2026, giving 31 calendar
days and 93 bread-day observations.  The manager's daily quantities are 3
Rustic Baguettes, 8 Demi Baguettes, and 1 Como Loaf.  For each bread $b$ and day
$t$, estimated sales are constructed as
\[
  \widehat D_{b,t}=Q^{\mathrm{current}}_b-W_{b,t},
\]
where $W_{b,t}$ is recorded waste.  A date absent from the transcribed waste
log is treated as having zero waste for all three products.  One suspected
duplicate set of September 30 entries is excluded.

The resulting dataset has no missing fields, duplicate bread-date keys, or
violations of the identity
\[
  \text{current order}-\text{waste}=\text{estimated sales}.
\]
Table~\ref{tab:audit} summarizes the observations used by the program.

\begin{table}[htbp]
\centering
\caption{Dataset audit by bread type.}
\label{tab:audit}
\begin{tabular}{@{}lrrrrr@{}}
\toprule
Bread & Days & Ordered & Estimated sales & Waste & Zero-waste days \\
\midrule
Rustic & 31 & 93  & 61  & 32 & 16 \\
Demi   & 31 & 248 & 213 & 35 & 19 \\
Como   & 31 & 31  & 20  & 11 & 20 \\
\midrule
Total  & 93 & 372 & 294 & 78 & 55 \\
\bottomrule
\end{tabular}
\end{table}

The available records do not show sellout time or attempted purchases after
inventory reached zero.  In addition, treating an unlisted date as zero waste
should be confirmed with the store's recordkeeping process.

\section{Optimization Rule}

For observations $D_1,\ldots,D_n$ in a group, the program sorts the values and
returns observation
\[
  D_{(\lceil 0.80n\rceil)}.
\]
This is the smallest integer order quantity whose empirical service level is
at least 80\%.  The fixed policy applies this rule to all 31 observations for
each bread.  The weekday policy applies it separately to groups of four or
five observations for each bread-weekday pair.

For a proposed quantity $Q$ and assumed demand $D$, the program evaluates
\[
  \text{waste}=\max(Q-D,0),\qquad
  \text{unmet demand}=\max(D-Q,0).
\]
Waste cost is waste multiplied by the corresponding wholesale unit cost.
Service level is the percentage of the 93 bread-day observations for which
$D\leq Q$.

\section{Baseline Results}

The baseline sets assumed demand equal to estimated sales, including on
zero-waste days.  The fixed optimized quantities are therefore identical to
the current quantities: Como 1, Demi 8, and Rustic 3.  The weekday result is
shown in Table~\ref{tab:schedule}.

\begin{table}[htbp]
\centering
\caption{Baseline weekday optimized quantities.}
\label{tab:schedule}
\begin{tabular}{@{}lrrr@{}}
\toprule
Weekday & Como & Demi & Rustic \\
\midrule
Monday    & 1 & 8 & 3 \\
Tuesday   & 1 & 8 & 3 \\
Wednesday & 1 & 8 & 2 \\
Thursday  & 1 & 8 & 3 \\
Friday    & 1 & 8 & 3 \\
Saturday  & 1 & 8 & 3 \\
Sunday    & 1 & 8 & 3 \\
\bottomrule
\end{tabular}
\end{table}

\begin{table}[htbp]
\centering
\caption{Baseline in-sample strategy comparison.}
\label{tab:baseline}
\begin{tabular}{@{}lrrrr@{}}
\toprule
Strategy & Wasted units & Waste cost & Service (\%) & Unmet units \\
\midrule
Current           & 78 & \$151.58 & 100.00 & 0 \\
Fixed optimized   & 78 & \$151.58 & 100.00 & 0 \\
Weekday optimized & 74 & \$141.62 & 100.00 & 0 \\
\bottomrule
\end{tabular}
\end{table}

The weekday policy reduces in-sample waste by 4 units, or
$4/78=5.13\%$, and reduces wholesale waste cost by \$9.96.  All four units are
Rustic Baguettes on Wednesdays.  These are in-sample calculations, not
out-of-sample performance estimates.

The 100\% baseline service level should not be interpreted as 100\% of actual
customer demand being met.  Estimated sales cannot exceed the historical
order quantity, so current orders cover every baseline observation by
construction.

\section{Exact Console Output}

The following output was produced by running \texttt{python}
\path{pasta_bread_cleaned_model.py} with Python 3.14, pandas 3.1.1, and
NumPy 2.5.4.

\begin{lstlisting}[style=output]
FIXED OPTIMAL QUANTITIES:
{'Como': 1, 'Demi': 8, 'Rustic': 3}

WEEKDAY OPTIMAL QUANTITIES:
           Como  Demi  Rustic
Monday        1     8       3
Tuesday       1     8       3
Wednesday     1     8       2
Thursday      1     8       3
Friday        1     8       3
Saturday      1     8       3
Sunday        1     8       3

STRATEGY COMPARISON:
                   wasted_units  waste_cost  service_level_pct  unmet_units
Current                    78.0      151.58              100.0          0.0
Fixed optimized            78.0      151.58              100.0          0.0
Weekday optimized          74.0      141.62              100.0          0.0
\end{lstlisting}

\section{Reproducibility and Limitations}

The baseline optimization and evaluation use the same observations.  The
reported 5.13\% reduction is therefore a historical fit and is not a validated
forecast.  Each weekday group contains only four or five observations.  With
four observations, the empirical 80th percentile is the maximum observation;
with five, it is the fourth ordered observation.  Weekday recommendations can
therefore change substantially when a single observation changes.

Additional waste records would improve sample size but would not by themselves
identify demand above the current stock level.  Stockout times, attempted
purchases after sellout, or observations collected under higher order
quantities would provide information about censored demand.

Future work will use sensitivity testing and Monte Carlo simulation to examine
how strongly the preliminary conclusions depend on the demand assumptions.

\section{Complete Python Source}

The complete analysis program is included directly from the repository below.
Keeping the listing linked to the source file prevents the appendix from
silently diverging from the executed code.

\lstinputlisting[style=code,language=Python]{pasta_bread_cleaned_model.py}

\end{document}
